\input{preamble.tex}

\title{Лабораторная работа № 5 \\ Конфигурирование VLAN}
\author{Сущенко Алина Николаевна \\ НПИбд-01-23}

\institute{Российский университет дружбы народов имени Патриса Лумумбы}
\date{2026}

\begin{document}

\frame{\titlepage}

\begin{frame}
\frametitle{Цель работы}
\begin{itemize}
    \item Получить основные навыки по настройке VLAN на коммутаторах сети с использованием протоколов VTP и Trunk.
    \item Научиться проверять изоляцию трафика между различными VLAN.
\end{itemize}
\end{frame}

\begin{frame}
\frametitle{Настройка Trunk-портов на msk-ansusenko-sw-01}
    \centering
    \includegraphics[width=0.8\textwidth]{../images/01.png}
    
    {\small Настройка trunk на портах g0/1 и g0/2}
\end{frame}

\begin{frame}
\frametitle{Настройка Trunk-портов на msk-ansusenko-sw-2}
    \centering
    \includegraphics[width=0.8\textwidth]{../images/02.png}
    
    {\small Настройка trunk на порту g0/2}
\end{frame}

\begin{frame}
\frametitle{Настройка Trunk-портов на msk-ansusenko-sw-3}
    \centering
    \includegraphics[width=0.8\textwidth]{../images/03.png}
    
    {\small Настройка trunk на порту g0/1}
\end{frame}

\begin{frame}
\frametitle{Настройка Trunk-портов на msk-ansusenko-sw-1}
    \centering
    \includegraphics[width=0.8\textwidth]{../images/04.png}
    
    {\small Настройка trunk на порту Fa0/24}
\end{frame}

\begin{frame}
\frametitle{Настройка VTP-сервера и создание VLAN}
    \centering
    \includegraphics[width=0.8\textwidth]{../images/05.png}
    
    {\small Настройка VTP-сервера на msk-ansusenko-sw-01}
\end{frame}

\begin{frame}
\frametitle{Настройка VTP-клиента на msk-ansusenko-sw-2}
    \centering
    \includegraphics[width=0.8\textwidth]{../images/06.png}
    
    {\small Настройка VTP-клиента}
\end{frame}

\begin{frame}
\frametitle{Настройка VTP-клиента на msk-ansusenko-sw-3}
    \centering
    \includegraphics[width=0.8\textwidth]{../images/07.png}
    
    {\small Настройка VTP-клиента}
\end{frame}

\begin{frame}
\frametitle{Настройка VTP-клиента и портов доступа на msk-ansusenko-sw-1}
    \centering
    \includegraphics[width=0.8\textwidth]{../images/08.png}
    
    {\small Настройка клиента и назначение портов в VLAN 101 и 104}
\end{frame}

\begin{frame}
\frametitle{Настройка VTP-клиента и диапазона портов на msk-ansusenko-sw-4}
    \centering
    \includegraphics[width=0.8\textwidth]{../images/09.png}
    
    {\small Настройка клиента и портов Fa0/1-5 в VLAN 101}
\end{frame}

\begin{frame}
\frametitle{Проверка VTP на сервере}
    \centering
    \includegraphics[width=0.8\textwidth]{../images/10.png}
    
    {\small show vtp status и show vlan brief на msk-ansusenko-sw-01}
\end{frame}

\begin{frame}
\frametitle{Проверка VTP на клиенте msk-ansusenko-sw-1}
    \centering
    \includegraphics[width=0.8\textwidth]{../images/11.png}
    
    {\small Статус VTP и список VLAN (Revision 13, порты Fa0/1 в VLAN 101, Fa0/20 в VLAN 104)}
\end{frame}

\begin{frame}
\frametitle{Проверка VTP на клиенте msk-ansusenko-sw-2}
    \centering
    \includegraphics[width=0.8\textwidth]{../images/12.png}
    
    {\small Статус VTP и список VLAN}
\end{frame}

\begin{frame}
\frametitle{Проверка VTP на клиенте msk-ansusenko-sw-3}
    \centering
    \includegraphics[width=0.8\textwidth]{../images/13.png}
    
    {\small Статус VTP и список VLAN (порт Fa0/11 в VLAN 103)}
\end{frame}

\begin{frame}
\frametitle{Начальное состояние VTP на msk-ansusenko-sw-4}
    \centering
    \includegraphics[width=0.8\textwidth]{../images/14.png}
    
    {\small Configuration Revision = 0, VLAN не синхронизированы}
\end{frame}

\begin{frame}
\frametitle{Проверка доступности VLAN 2 с сервера}
    \centering
    \includegraphics[width=0.8\textwidth]{../images/15.png}
    
    {\small Пинг до коммутаторов 10.128.1.3, 10.128.1.4, 10.128.1.6}
\end{frame}

\begin{frame}
\frametitle{Проверка доступности VLAN 2 с msk-ansusenko-sw-4 после синхронизации}
    \centering
    \includegraphics[width=0.8\textwidth]{../images/16.png}
    
    {\small Успешные пинги подтверждают синхронизацию VTP}
\end{frame}

\begin{frame}
\frametitle{Проверка доступности VLAN 2 с msk-ansusenko-sw-1}
    \centering
    \includegraphics[width=0.8\textwidth]{../images/17.png}
    
    {\small Пинг до коммутаторов 10.128.1.3 и 10.128.1.4}
\end{frame}

\begin{frame}
\frametitle{Проверка доступности устройств в разных VLAN}
    \centering
    \includegraphics[width=0.8\textwidth]{../images/18.png}
    
    {\small Пинг на 10.128.101.10 (собственный адрес) и 10.128.103.10 (другой VLAN)}
\end{frame}

\begin{frame}
\frametitle{Проверка доступности устройств в разных VLAN}
    \centering
    \includegraphics[width=0.8\textwidth]{../images/19.png}
    
    {\small Движение пакетов в режиме симуляции}
\end{frame}

\begin{frame}
\frametitle{События симуляции в Packet Tracer}
    \centering
    \includegraphics[width=0.8\textwidth]{../images/20.png}
    
    {\small ICMP, STP и CDP пакеты в режиме симуляции}
\end{frame}

\begin{frame}
\frametitle{Сводная информация по настройкам VLAN}
В процессе выполнения работы были настроены следующие VLAN и назначены порты
\end{frame}

\begin{frame}
\frametitle{Контрольные вопросы}

\begin{block}{Вопрос 1: Какая команда используется для просмотра списка VLAN на сетевом устройстве?}
Для просмотра списка VLAN используется команда:
\begin{itemize}
    \item \texttt{show vlan brief} - отображает краткую информацию о VLAN и назначенных портах
    \item \texttt{show vlan} - подробная информация о VLAN
    \item \texttt{show vlan id <номер>} - информация о конкретном VLAN
\end{itemize}
\end{block}
\end{frame}

\begin{frame}
\frametitle{Контрольные вопросы}

\begin{block}{Вопрос 2: Охарактеризуйте VLAN Trunking Protocol (VTP). Приведите перечень команд с пояснениями для настройки и просмотра информации о VLAN.}
VTP (VLAN Trunking Protocol) - протокол Cisco для централизованного управления VLAN в домене. Позволяет синхронизировать информацию о VLAN между коммутаторами.

\textbf{Команды настройки VTP:}
\begin{itemize}
    \item \texttt{vtp mode \{server | client | transparent\}} - задать режим работы
    \item \texttt{vtp domain <имя>} - задать домен VTP
    \item \texttt{vtp password <пароль>} - установить пароль для аутентификации
\end{itemize}

\textbf{Команды просмотра информации:}
\begin{itemize}
    \item \texttt{show vtp status} - просмотр статуса VTP
    \item \texttt{show vlan brief} - просмотр списка VLAN
    \item \texttt{show interfaces trunk} - просмотр транковых портов
\end{itemize}
\end{block}
\end{frame}

\begin{frame}
\frametitle{Контрольные вопросы}

\begin{block}{Вопрос 3: Охарактеризуйте Internet Control Message Protocol (ICMP). Опишите формат пакета ICMP.}
ICMP (Internet Control Message Protocol) - протокол сетевого уровня, используемый для диагностики и сообщений об ошибках. Основные функции: проверка доступности (ping), трассировка маршрута (traceroute), сообщения об ошибках.

\textbf{Формат пакета ICMP:}
\begin{itemize}
    \item \textbf{Тип (8 бит)} - определяет тип сообщения (8 — echo request, 0 — echo reply)
    \item \textbf{Код (8 бит)} - уточняет тип сообщения
    \item \textbf{Контрольная сумма (16 бит)} - проверка целостности
    \item и другие поля в зависимости от типа сообщения
\end{itemize}
\end{block}
\end{frame}

\begin{frame}
\frametitle{Контрольные вопросы}

\begin{block}{Вопрос 4: Охарактеризуйте Address Resolution Protocol (ARP). Опишите формат пакета ARP.}
ARP (Address Resolution Protocol) - протокол для сопоставления IP-адреса с MAC-адресом в локальной сети. Используется для определения MAC-адреса получателя по известному IP-адресу.

\textbf{Формат пакета ARP:}
\begin{itemize}
    \item \textbf{Тип оборудования (16 бит)} - 1 для Ethernet
    \item \textbf{Тип протокола (16 бит)} - 0x0800 для IPv4
    \item \textbf{Длина оборудования (8 бит)} - 6 для MAC-адреса
    \item \textbf{Длина протокола (8 бит)} - 4 для IPv4
    \item \textbf{Операция (16 бит)} - 1 для запроса, 2 для ответа
    \item \textbf{MAC отправителя (48 бит)}
    \item \textbf{IP отправителя (32 бит)}
    \item \textbf{MAC цели (48 бит)}
    \item \textbf{IP цели (32 бит)}
\end{itemize}
\end{block}
\end{frame}

\begin{frame}
\frametitle{Контрольные вопросы}

\begin{block}{Вопрос 5: Что такое MAC-адрес? Какова его структура?}
MAC-адрес (Media Access Control) - уникальный идентификатор сетевого интерфейса на канальном уровне модели OSI. Используется для адресации устройств в пределах локальной сети.

\textbf{Структура MAC-адреса (6 байт = 48 бит):}
\begin{itemize}
    \item \textbf{Первые 3 байта (24 бита)} - OUI (Organizationally Unique Identifier), идентификатор производителя оборудования
    \item \textbf{Последние 3 байта (24 бита)} - уникальный номер устройства, назначаемый производителем
\end{itemize}

\textbf{Форматы записи:}
\begin{itemize}
    \item \texttt{00:11:22:AA:BB:CC} (двоеточие)
    \item \texttt{00-11-22-AA-BB-CC} (дефис)
    \item \texttt{0011.22AA.BBCC} (точка)
\end{itemize}
\end{block}
\end{frame}

\begin{frame}
\frametitle{Выводы}
\begin{itemize}
    \item В результате выполнения работы были настроены Trunk-порты на всех коммутаторах в соответствии с топологией сети.
    \item Настроен VTP-домен с сервером и клиентами для централизованного управления VLAN.
    \item Проверена изоляция трафика между различными VLAN.
\end{itemize}
\end{frame}

\end{document}
